\documentclass[10pt,a4paper]{amsart}
\usepackage{amsmath,amssymb,amsthm,graphicx}
\usepackage{verbatim}
\newenvironment{reference}{\comment}{\endcomment}
%\newenvironment{reference}{}{}
\newenvironment{slogan}{\comment}{\endcomment}
\newenvironment{history}{\comment}{\endcomment}

% For commutative diagrams we use Xy-pic
\usepackage[all]{xy}

% We use 2cell for 2-commutative diagrams.
\xyoption{2cell}
% Theorem environments.
%
\theoremstyle{plain}
\newtheorem{theorem}[subsection]{Theorem}
\newtheorem{proposition}[subsection]{Proposition}
\newtheorem{lemma}[subsection]{Lemma}

\theoremstyle{definition}
\newtheorem{definition}[subsection]{Definition}
\newtheorem{example}[subsection]{Example}
\newtheorem{exercise}[subsection]{Exercise}
\newtheorem{situation}[subsection]{Situation}

\theoremstyle{remark}
\newtheorem{remark}[subsection]{Remark}
\newtheorem{remarks}[subsection]{Remarks}

\numberwithin{equation}{subsection}

% Macros
%
\def\lim{\mathop{\mathrm{lim}}\nolimits}
\def\colim{\mathop{\mathrm{colim}}\nolimits}
\def\Spec{\mathop{\mathrm{Spec}}}
\def\Hom{\mathop{\mathrm{Hom}}\nolimits}
\def\Ext{\mathop{\mathrm{Ext}}\nolimits}
\def\SheafHom{\mathop{\mathcal{H}\!\mathit{om}}\nolimits}
\def\SheafExt{\mathop{\mathcal{E}\!\mathit{xt}}\nolimits}
\def\Sch{\mathit{Sch}}
\def\Mor{\mathop{\mathrm{Mor}}\nolimits}
\def\Ob{\mathop{\mathrm{Ob}}\nolimits}
\def\Sh{\mathop{\mathit{Sh}}\nolimits}
\def\NL{\mathop{N\!L}\nolimits}
\def\CH{\mathop{\mathrm{CH}}\nolimits}
\def\proetale{{pro\text{-}\acute{e}tale}}
\def\etale{{\acute{e}tale}}
\def\QCoh{\mathit{QCoh}}
\def\Ker{\mathop{\mathrm{Ker}}}
\def\Im{\mathop{\mathrm{Im}}}
\def\Coker{\mathop{\mathrm{Coker}}}
\def\Coim{\mathop{\mathrm{Coim}}}

% Boxtimes
%
\DeclareMathSymbol{\boxtimes}{\mathbin}{AMSa}{"02}

%
% Macros for moduli stacks/spaces
%
\def\QCohstack{\mathcal{QC}\!\mathit{oh}}
\def\Cohstack{\mathcal{C}\!\mathit{oh}}
\def\Spacesstack{\mathcal{S}\!\mathit{paces}}
\def\Quotfunctor{\mathrm{Quot}}
\def\Hilbfunctor{\mathrm{Hilb}}
\def\Curvesstack{\mathcal{C}\!\mathit{urves}}
\def\Polarizedstack{\mathcal{P}\!\mathit{olarized}}
\def\Complexesstack{\mathcal{C}\!\mathit{omplexes}}
% \Pic is the operator that assigns to X its picard group, usage \Pic(X)
% \Picardstack_{X/B} denotes the Picard stack of X over B
% \Picardfunctor_{X/B} denotes the Picard functor of X over B
\def\Pic{\mathop{\mathrm{Pic}}\nolimits}
\def\Picardstack{\mathcal{P}\!\mathit{ic}}
\def\Picardfunctor{\mathrm{Pic}}
\def\Deformationcategory{\mathcal{D}\!\mathit{ef}}
\usepackage{cleveref}
\crefname{originalreference}{original reference}{original references}
\Crefname{originalreference}{Original reference}{Original references}
\makeatletter
\newlabel{lemma-injective-colimit-direct-sum-pushforwards-artin}{{09TJ}{source}{}{sourcepointer-1}{}}
\newlabel{lemma-injective-colimit-direct-sum-pushforwards-artin@cref}{{[originalreference][0][]09TJ}{[1][0][]source}{}{}{}}
\makeatother
\crefname{originallemma}{original lemma}{original lemmas}
\Crefname{originallemma}{Original lemma}{Original lemmas}
\makeatletter
\newlabel{sites-cohomology-lemma-apply-Leray}{{0733}{source}{}{sourcepointer-2}{}}
\newlabel{sites-cohomology-lemma-apply-Leray@cref}{{[originallemma][0][]0733}{[1][0][]source}{}{}{}}
\makeatother
\crefname{originallemma}{original lemma}{original lemmas}
\Crefname{originallemma}{Original lemma}{Original lemmas}
\makeatletter
\newlabel{spaces-cohomology-lemma-affine-base-change}{{07U8}{source}{}{sourcepointer-3}{}}
\newlabel{spaces-cohomology-lemma-affine-base-change@cref}{{[originallemma][0][]07U8}{[1][0][]source}{}{}{}}
\makeatother
\crefname{originallemma}{original lemma}{original lemmas}
\Crefname{originallemma}{Original lemma}{Original lemmas}
\makeatletter
\newlabel{spaces-cohomology-lemma-colimits}{{073E}{source}{}{sourcepointer-4}{}}
\newlabel{spaces-cohomology-lemma-colimits@cref}{{[originallemma][0][]073E}{[1][0][]source}{}{}{}}
\makeatother
\crefname{originallemma}{original lemma}{original lemmas}
\Crefname{originallemma}{Original lemma}{Original lemmas}
\makeatletter
\newlabel{spaces-morphisms-lemma-Artinian-affine}{{09TF}{source}{}{sourcepointer-5}{}}
\newlabel{spaces-morphisms-lemma-Artinian-affine@cref}{{[originallemma][0][]09TF}{[1][0][]source}{}{}{}}
\makeatother
\title{Injective quasi-coherent sheaves are acyclic etale-locally and have vanishing higher direct images}
\author{The Stacks Project Authors}
\date{stacks-project:spaces-perfect.tex:local-snapshot-20261012; GNU FDL 1.2 or later}
\begin{document}
\maketitle
\noindent\textit{Editorial note.} All delivered mathematical text is the author's own, copied from the pinned Stacks Project chapter file spaces-perfect.tex; no statement or proof step was rewritten, reordered or removed.\par\medskip
\noindent\textit{Editorial note.} The counted proof is the author's complete argument: triage\_override: the author's '(details omitted)' sits on 'V x\_Y X is qcqs etale over X', justified by etale base-change stability plus quasi-compactness of morphisms from Noetherian spaces; part (1) additionally leans on the delivered L48106 structure lemma\par\medskip
\noindent\textit{Editorial note.} The label lemma-injective-pushforward names the counted statement, and the proof environment immediately following it in the source is delivered verbatim as its proof.\par\medskip
\noindent\textit{Editorial note.} The counted statement is carried by the author's own lemma environment in the chapter source.\par\medskip
\noindent\textit{Allowed standard prerequisites (editorial).}\begin{itemize}
\item Lemma lemma-injective-colimit-direct-sum-pushforwards-artin (delivered, L48106): every injective in QCoh(O\_X) is a direct summand of a colimit of pushforwards from affine closed-ish subschemes over a Noetherian scheme... (the structure lemma used to decompose J)
\item Cohomology of Spaces, Lemma spaces-cohomology-lemma-colimits: H\textasciicircum{}p commutes with filtered colimits on quasi-compact quasi-separated spaces
\item Cohomology of Spaces, Lemma spaces-cohomology-lemma-affine-base-change: affine base change identification of the pullback
\item Cohomology on Sites, Lemma sites-cohomology-lemma-apply-Leray: Leray spectral sequence
\end{itemize}\medskip
\ifdefined\hypertarget\hypertarget{sourcepointer-1}{}\fi
\noindent\textit{Source reference.} Original reference 09TJ: sources/stacks-spaces-perfect/pinned-original-source.tex, source lines 2868--2868. Prerequisite pointer invoked at 2932; it names a standard fact of the same chapter that the counted argument uses, and every step of the counted proof built on it is retained verbatim in the delivered body.\par\medskip
\ifdefined\hypertarget\hypertarget{sourcepointer-2}{}\fi
\noindent\textit{Source reference.} Original lemma 0733: sources/stacks-chow/tags.txt, source lines 9047--9047. Prerequisite pointer invoked at 2952; it names a standard fact of a sibling chapter that the counted argument uses, and every step of the counted proof built on it is retained verbatim in the delivered body.\par\medskip
\ifdefined\hypertarget\hypertarget{sourcepointer-3}{}\fi
\noindent\textit{Source reference.} Original lemma 07U8: sources/stacks-chow/tags.txt, source lines 9963--9963. Prerequisite pointer invoked at 2947; it names a standard fact of a sibling chapter that the counted argument uses, and every step of the counted proof built on it is retained verbatim in the delivered body.\par\medskip
\ifdefined\hypertarget\hypertarget{sourcepointer-4}{}\fi
\noindent\textit{Source reference.} Original lemma 073E: sources/stacks-chow/tags.txt, source lines 9058--9058. Prerequisite pointer invoked at 2937; it names a standard fact of a sibling chapter that the counted argument uses, and every step of the counted proof built on it is retained verbatim in the delivered body.\par\medskip
\ifdefined\hypertarget\hypertarget{sourcepointer-5}{}\fi
\noindent\textit{Source reference.} Original lemma 09TF: sources/stacks-chow/tags.txt, source lines 12388--12388. Prerequisite pointer invoked at 2939; it names a standard fact of a sibling chapter that the counted argument uses, and every step of the counted proof built on it is retained verbatim in the delivered body.\par\medskip
\begin{lemma}
\label{lemma-injective-pushforward}
Let $S$ be a scheme. Let $X$ be a Noetherian algebraic space over $S$. If
$\mathcal{J}$ is an injective object of $\QCoh(\mathcal{O}_X)$,
then
\begin{enumerate}
\item $H^p(U, \mathcal{J}|_U) = 0$ for $p > 0$ and for
every quasi-compact and quasi-separated algebraic space $U$ \'etale over $X$,
\item for any morphism $f : X \to Y$ of algebraic spaces over $S$
with $Y$ quasi-separated we have $R^pf_*\mathcal{J} = 0$ for $p > 0$.
\end{enumerate}
\end{lemma}

\begin{proof}
Proof of (1). Write $\mathcal{J}$ as a direct summand of
$\colim \mathcal{F}_i$ with $\mathcal{F}_i = (Z_i \to X)_*\mathcal{G}_i$
as in Lemma \ref{lemma-injective-colimit-direct-sum-pushforwards-artin}.
It is clear that it suffices to prove the vanishing for
$\colim \mathcal{F}_i$. Since pullback commutes with colimits
and since $U$ is quasi-compact and quasi-separated, it suffices
to prove $H^p(U, \mathcal{F}_i|_U) = 0$ for $p > 0$, see
Cohomology of Spaces, Lemma \ref{spaces-cohomology-lemma-colimits}.
Observe that $Z_i \to X$ is an affine morphism, see
Morphisms of Spaces, Lemma \ref{spaces-morphisms-lemma-Artinian-affine}.
Thus
$$
\mathcal{F}_i|_U = (Z_i \times_X U \to U)_*\mathcal{G}'_i =
R(Z_i \times_X U \to U)_*\mathcal{G}'_i
$$
where $\mathcal{G}'_i$ is the pullback of $\mathcal{G}_i$
to $Z_i \times_X U$, see
Cohomology of Spaces, Lemma \ref{spaces-cohomology-lemma-affine-base-change}.
Since $Z_i \times_X U$ is affine we conclude that
$\mathcal{G}'_i$ has no higher cohomology on $Z_i \times_X U$.
By the Leray spectral sequence we conclude the same
thing is true for $\mathcal{F}_i|_U$ (Cohomology on Sites,
Lemma \ref{sites-cohomology-lemma-apply-Leray}).

\medskip\noindent
Proof of (2). Let $f : X \to Y$ be a morphism of algebraic spaces
over $S$. Let $V \to Y$ be an \'etale morphism with $V$ affine.
Then $V \times_Y X \to X$ is an \'etale morphism and
$V \times_Y X$ is a quasi-compact and quasi-separated algebraic
space \'etale over $X$ (details omitted). Hence
$H^p(V \times_Y X, \mathcal{J})$ is zero by part (1).
Since $R^pf_*\mathcal{J}$ is the sheaf associated to the presheaf
$V \mapsto H^p(V \times_Y X, \mathcal{J})$ the result is proved.
\end{proof}

\end{document}
